\documentclass[]{article}
\makeatletter
\renewcommand\paragraph{\@startsection{paragraph}{4}{\z@}%
            {-2.5ex\@plus -1ex \@minus -.25ex}%
            {1.25ex \@plus .25ex}%
            {\normalfont\normalsize\bfseries}}
\makeatother
\setcounter{secnumdepth}{4}
\usepackage{amsmath}
%\usepackage{mathtools}
\usepackage{graphicx}
\usepackage{slashed}
\usepackage{lineno}
\usepackage{latexsym}
\usepackage{subfigure}
\usepackage{amssymb}
\newtheorem{thm}{Theorem}[section]
\newtheorem{cor}[thm]{Corollary}
\newtheorem{lem}[thm]{Lemma}
\usepackage[numbers,sort]{natbib}

\newcommand{\bb}{\begin{equation}}
\newcommand{\ee}{\end{equation}}
\newtheorem{defin}{Definition}
\usepackage{multirow}
\usepackage{ctable}
\usepackage{bm}
\usepackage{enumerate}
\newcommand{\D}[2]{\frac{\partial #1}{\partial #2}}
\newcommand{\DD}[2]{\frac{\partial^2 #1}{\partial #2^2}}
\newcommand{\rd}{\text{ d}}
\usepackage{framed}
\newcommand{\see}[1]{(see Figure \ref{#1})}
\newcommand{\fig}[1]{Figure \ref{#1}}
\newcommand{\figs}[2]{figures \ref{#1} and \ref{#2}}
\newcommand{\sect}[1]{Section \ref{#1}}
\newcommand{\app}[1]{Appendix \ref{#1}}
\newcommand{\chap}[1]{Chapter \ref{#1}}
\newcommand{\eqn}[1]{equation \eqref{#1}}
\newcommand{\eqns}[2]{equations \eqref{#1} and \eqref{#2}}
\newcommand{\eqnto}[2]{equations \eqref{#1}-\eqref{#2}}
%\usepackage{authblk}
\usepackage{url}
\usepackage{soul}
\newcommand{\eg}{\emph{e.g.} }
\newcommand{\bn}{\bm{n}}
\newcommand{\bu}{\bm{u}}
\newcommand{\ie}{\emph{i.e.} }
\newcommand{\Chapter}[1]{\chapter{#1}\label{#1}}
\newcommand{\Section}[1]{\section{#1}\label{#1}}
\newcommand{\Subsection}[1]{\subsection{#1}\label{#1}}
\newcommand{\Subsubsection}[1]{\subsubsection{#1}\label{#1}}
\newcommand{\Appendix}[1]{\appendix{#1}\label{#1}}
\usepackage[margin=3cm,centering]{geometry}
\usepackage[geometry]{ifsym}
\makeatletter
\newcommand\restr[2]{{% we make the whole thing an ordinary symbol
  \left.\kern-\nulldelimiterspace % automatically resize the bar with \right
  #1 % the function
  \vphantom{\big|} % pretend it's a little taller at normal size
  \right|_{#2} % this is the delimiter
  }}
\def\url@leostyle{%
  \@ifundefined{selectfont}{\def\UrlFont{\sf}}{\def\UrlFont{\small\ttfamily}}}
\makeatother
\urlstyle{leo}
\usepackage{multirow}
\usepackage{blkarray}
\usepackage{soul}
\usepackage{framed}
\usepackage{color}
\usepackage{setspace}
\newcommand{\ttttp}{.24\textwidth}
\newcommand{\tttp}{.32\textwidth}
\newcommand{\ttp}{.45\textwidth}
\newcommand{\tbo}{.6\textwidth}
 \usepackage[T1]{fontenc}
\usepackage[utf8]{inputenc}
\usepackage{authblk}
 \renewcommand{\l}{\left(}
\renewcommand{\r}{\right)}

%\begin{figure}[h!!!tb]
%\centering
%\subfigure[\label{Godzilla}]{\includegraphics[height=0.35\textwidth]{./Pictures/Godzilla_final_bw.png}}
%\subfigure[\label{Jaeger}]{\includegraphics[height=0.35\textwidth]{./Pictures/Jaeger_finish_bw.png}}
%\caption{\label{Monsters} The two types of monster we are going to consider are: (a) the naturally occurring Kaijus and (b) the man-made Jaegers.}
%\end{figure}
\title{Problem sheet 1}
\author{Thomas E. Woolley}
\begin{document}
\maketitle
\begin{enumerate}
\item \textbf{Carbon dating}
\item \textbf{Drag racing}
\item \textbf{Newton's law of cooling}
You can find the temperature inside your refrigerator without putting
a thermometer inside. Take a can of soda from the refrigerator, let
it warm for half an hour, then record its temperature. Let it warm
for another half an hour and record its temperature again. Suppose
that the readings are
T
(1
=
2) = 45
�
F
and
T
(1) = 55
�
F
. Assuming
that the room temperature is
70
�
F
, what is the temperature inside the
refrigerator?
Solution:
Taking the time one half hour after the soda was removed from
the refrigerator to be the ``zero time'' (and stating the given information in
an appropriate way), we have the boundary value problem
dT
dt
=
k
(70




\item \textbf{Motion under gravity}
Motion under gravity, throwing a ball straight up 

In this case we want to consider the vertical displacement of a projectile. We define $u$ to be the more standard coordinate variable $y$. We then know from physics that the acceleration (second) of the particle is equal to the gravitational constant $g=9.8$m/s$^2$.

\item \textbf{Logistic equation}
\begin{align}
\frac{\rd u}{\rd t}&=rcu\l 1-\frac{u}{c}\r,\nonumber\\
\Rightarrow\int_0^T\frac{\rd u}{u(1-u/c)}&=\int_0^Trc \rd t,\nonumber\\
\Rightarrow\int_0^T \frac{1}{u}+\frac{1/c}{1-u/c}  \rd u&=rcT,\nonumber\\
\Rightarrow\left[\ln(u)-\ln\l 1- \frac{u}{c} \r \right]^T_0&=rcT,\nonumber\\
\Rightarrow\ln\l\frac{u}{1- \frac{u}{c}} \r-\ln\l\frac{u_0}{1- \frac{u_0}{c}} \r&=rcT,\nonumber\\
\Rightarrow u(T)=\frac{c}{1+\frac{c-u_0}{u_0}\exp\l{-rcT}\r}
\end{align}

\item \textbf{Bacteria growth}
\dot{u}=v
\dot{v}=-uv

\item \textbf{Time shift in ODEs, arbitrary start point}
\item \textbf{Taylor series}
\item \textbf{Wobbling doors} Damped oscillator.

\item \textbf{Two body problem}
We can combine the Second Law of Motion and the Law of Gravitation in order to predict the positions of planets interacting through their gravitational fields \see{Two_body}. Consider two planets with the same mass, $m$, and positions $\bm{r}(t)_1$ and $\bm{r}(t)_2$, respectively. Further, we note that since we are dealing with acceleration (a second-order equation) we will need to specify two initial conditions, the position and velocity. Let the initial positions be $\bm{r}(0)_1=\bm{r}_{10}$ and $\bm{r}(0)_2=\bm{r}_{20}$ and the initial velocities $\dot{\bm{r}(0)_1}=\bm{v}_{10}$ and $\dot{\bm{r}(0)_2}=\bm{v}_{20}$. The governing equations are
\begin{align}
&m\ddot{\bm{r}_1}=-G\frac{m^2}{|\bm{r}_1-\bm{r}_2|^3}\l\bm{r}_1-\bm{r}_2\r,\label{g1}\\
&m\ddot{\bm{r}_2}=-G\frac{m^2}{|\bm{r}_2-\bm{r}_1|^3}\l\bm{r}_2-\bm{r}_1\r\label{g2}.
\end{align}
where we remember that this is a vector equation, $\bm{r}_i=(x_i,y_i)$, so there are actually four ODEs here, rather than two.

We notice that if we add the equations together the right-hand sides cancel out,
\bb
\ddot{\bm{r}_1}+\ddot{\bm{r}_2}=0.
\ee
Integrating this equation we find that the velocity of $\bm{r}_1+\bm{r}_2$ is a conserved quantity and that
\bb
\bm{r}_1+\bm{r}_2=(\bm{r}_{10}+\bm{r}_{20})+(\bm{v}_{10}+\bm{v}_{20})t.\label{COM}
\ee
Points to note:
\begin{itemize}
\item the defined quantity $\bm{R}=(\bm{r}_1+\bm{r}_2)/2$ is the centre of mass of the system (because $m_1=m_2$). Thus the motion of the two planets is constrained such that their centre of mass travels at a constant speed throughout space;
\item using \eqn{COM} we could decouple \eqns{g1}{g2} into two separate non-autonomous ODEs;
\item it turns out to be more fruitful to consider the quantity $\bm{r}_1-\bm{r}_2$, which, if written in polar coordinates, leads to an ODE with a closed form solution. We will not do this here, but we will consider the simulated results.
\item when solving \eqns{g1}{g2} numerically we could separate each equation into its Cartesian components and reduce the second-order equation to two first-order equations. Namely, we introduce $(v_{ix},v_{iy})=(\dot{x}_i,\dot{y}_i)$. Thus, we would have a system of eight ODEs to solve, with variables $(x_1,y_1,v_{1x},v_{1y},x_2,y_2,v_{2x},v_{2y})$. However, since $x_i$ and $y_i$ are perpendicular Cartesian coordinates it turns out to be a good idea to use complex numbers. Specifically, instead of writing two ODEs for each of $x_i$ and $y_i$ we can simply solve one ODE in terms of the complex quantity $\bm{r}_i=x_i+Iy_i$, which can be handled by numerical solvers. Thus, we simplify the numerical solution from eight equations to four equations.
\end{itemize}

\end{enumerate}
\end{document}

